% Part: normal-modal-logic
% Chapter: sequent-calculus
% Section: introduction

\documentclass[../../../include/open-logic-section]{subfiles}

\begin{document}

\olfileid{nml}{seq}{int}

\olsection{Pambuka}

Kalkulus sekuen kanggo logika proposisional bisa diambakake nganggo
aturan tambahan kang nangani
\iftag{prvBox}{$\Box$\iftag{prvDiamond}{ lan~}{}}{}%
\iftag{prvDiamond}{$\Diamond$}{}. Contone, kanggo~$\Log{K}$,
kita nduweni \Log{LK} kanthi tambahan:
\[
  \iftag{prvBox}{\iftag{prvDiamond}{% <> and [] primitive
    \Axiom$\Gamma \fCenter \Delta, !A$
    \RightLabel{$\Box$}
    \UnaryInf$\Box\Gamma \fCenter \Diamond\Delta, \Box!A$
    \DisplayProof  
    \qquad
    \Axiom$!A, \Gamma \fCenter \Delta$
    \RightLabel{$\Diamond$}
    \UnaryInf$\Diamond!A, \Box\Gamma \fCenter \Diamond\Delta$
    \DisplayProof
}{% only Box primitive
\Axiom$\Gamma \fCenter !A$
\RightLabel{$\Box$}
\UnaryInf$\Box\Gamma \fCenter \Box!A$
\DisplayProof
}}{% only <> primitive
  \Axiom$!A \fCenter \Delta$
  \RightLabel{$\Diamond$}
  \UnaryInf$\Diamond!A \fCenter \Diamond\Delta$
  \DisplayProof
}
\]
Kanggo perluasan~$\Log{K}$, aturan tambahan uga kudu ditambahake.

Ora saben logika modal nduweni kalkulus sekuen kang kaya mangkono.
Malah $\Log{S5}$, kang semantike prasaja (bisa ditegesi tanpa migunakake
relasi aksesibilitas babar pisan), durung dikenal nduweni kalkulus
sekuen turunan saka~$\Log{LK}$ kang lengkap tanpa aturan~\Cut.
Nanging, logika iki nduweni kalkulus \emph{hipersekuen} kang lengkap
lan tanpa aturan cut.

\end{document}
